\documentclass[10pt,a4paper]{amsart}
\usepackage{amsmath,amssymb,amsthm,graphicx}
\usepackage{verbatim}
\newenvironment{reference}{\comment}{\endcomment}
%\newenvironment{reference}{}{}
\newenvironment{slogan}{\comment}{\endcomment}
\newenvironment{history}{\comment}{\endcomment}

% For commutative diagrams we use Xy-pic
\usepackage[all]{xy}

% We use 2cell for 2-commutative diagrams.
\xyoption{2cell}
% Theorem environments.
%
\theoremstyle{plain}
\newtheorem{theorem}[subsection]{Theorem}
\newtheorem{proposition}[subsection]{Proposition}
\newtheorem{lemma}[subsection]{Lemma}

\theoremstyle{definition}
\newtheorem{definition}[subsection]{Definition}
\newtheorem{example}[subsection]{Example}
\newtheorem{exercise}[subsection]{Exercise}
\newtheorem{situation}[subsection]{Situation}

\theoremstyle{remark}
\newtheorem{remark}[subsection]{Remark}
\newtheorem{remarks}[subsection]{Remarks}

\numberwithin{equation}{subsection}

% Macros
%
\def\lim{\mathop{\mathrm{lim}}\nolimits}
\def\colim{\mathop{\mathrm{colim}}\nolimits}
\def\Spec{\mathop{\mathrm{Spec}}}
\def\Hom{\mathop{\mathrm{Hom}}\nolimits}
\def\Ext{\mathop{\mathrm{Ext}}\nolimits}
\def\SheafHom{\mathop{\mathcal{H}\!\mathit{om}}\nolimits}
\def\SheafExt{\mathop{\mathcal{E}\!\mathit{xt}}\nolimits}
\def\Sch{\mathit{Sch}}
\def\Mor{\mathop{\mathrm{Mor}}\nolimits}
\def\Ob{\mathop{\mathrm{Ob}}\nolimits}
\def\Sh{\mathop{\mathit{Sh}}\nolimits}
\def\NL{\mathop{N\!L}\nolimits}
\def\CH{\mathop{\mathrm{CH}}\nolimits}
\def\proetale{{pro\text{-}\acute{e}tale}}
\def\etale{{\acute{e}tale}}
\def\QCoh{\mathit{QCoh}}
\def\Ker{\mathop{\mathrm{Ker}}}
\def\Im{\mathop{\mathrm{Im}}}
\def\Coker{\mathop{\mathrm{Coker}}}
\def\Coim{\mathop{\mathrm{Coim}}}

% Boxtimes
%
\DeclareMathSymbol{\boxtimes}{\mathbin}{AMSa}{"02}

%
% Macros for moduli stacks/spaces
%
\def\QCohstack{\mathcal{QC}\!\mathit{oh}}
\def\Cohstack{\mathcal{C}\!\mathit{oh}}
\def\Spacesstack{\mathcal{S}\!\mathit{paces}}
\def\Quotfunctor{\mathrm{Quot}}
\def\Hilbfunctor{\mathrm{Hilb}}
\def\Curvesstack{\mathcal{C}\!\mathit{urves}}
\def\Polarizedstack{\mathcal{P}\!\mathit{olarized}}
\def\Complexesstack{\mathcal{C}\!\mathit{omplexes}}
% \Pic is the operator that assigns to X its picard group, usage \Pic(X)
% \Picardstack_{X/B} denotes the Picard stack of X over B
% \Picardfunctor_{X/B} denotes the Picard functor of X over B
\def\Pic{\mathop{\mathrm{Pic}}\nolimits}
\def\Picardstack{\mathcal{P}\!\mathit{ic}}
\def\Picardfunctor{\mathrm{Pic}}
\def\Deformationcategory{\mathcal{D}\!\mathit{ef}}
\usepackage{cleveref}
\crefname{originalreference}{original reference}{original references}
\Crefname{originalreference}{Original reference}{Original references}
\makeatletter
\newlabel{lemma-flat-proper-perfect-direct-image-general}{{0B91}{source}{}{sourcepointer-1}{}}
\newlabel{lemma-flat-proper-perfect-direct-image-general@cref}{{[originalreference][0][]0B91}{[1][0][]source}{}{}{}}
\makeatother
\crefname{originalreference}{original reference}{original references}
\Crefname{originalreference}{Original reference}{Original references}
\makeatletter
\newlabel{lemma-jump-loci}{{0BDI}{source}{}{sourcepointer-2}{}}
\newlabel{lemma-jump-loci@cref}{{[originalreference][0][]0BDI}{[1][0][]source}{}{}{}}
\makeatother
\crefname{originalreference}{original reference}{original references}
\Crefname{originalreference}{Original reference}{Original references}
\makeatletter
\newlabel{lemma-locally-closed-where-H0-locally-free}{{0BDL}{source}{}{sourcepointer-3}{}}
\newlabel{lemma-locally-closed-where-H0-locally-free@cref}{{[originalreference][0][]0BDL}{[1][0][]source}{}{}{}}
\makeatother
\crefname{originallemma}{original lemma}{original lemmas}
\Crefname{originallemma}{Original lemma}{Original lemmas}
\makeatletter
\newlabel{more-algebra-lemma-better-cut-complex-in-two}{{0A1U}{source}{}{sourcepointer-4}{}}
\newlabel{more-algebra-lemma-better-cut-complex-in-two@cref}{{[originallemma][0][]0A1U}{[1][0][]source}{}{}{}}
\makeatother
\title{Proper flat morphisms with trivial degree-zero function field cohomology}
\author{The Stacks Project Authors}
\date{stacks-project:spaces-perfect.tex:local-snapshot-20261012; GNU FDL 1.2 or later}
\begin{document}
\maketitle
\noindent\textit{Editorial note.} All delivered mathematical text is the author's own, copied from the pinned Stacks Project chapter file spaces-perfect.tex; no statement or proof step was rewritten, reordered or removed.\par\medskip
\noindent\textit{Editorial note.} The counted proof is the author's complete argument: both clauses are proved in full, the scheme case being handled through the cited jump-locus and cut-complex lemmas\par\medskip
\noindent\textit{Editorial note.} The label lemma-proper-flat-h0 names the counted statement, and the proof environment immediately following it in the source is delivered verbatim as its proof.\par\medskip
\noindent\textit{Editorial note.} The counted statement is carried by the author's own lemma environment in the chapter source.\par\medskip
\noindent\textit{Allowed standard prerequisites (editorial).}\begin{itemize}
\item cohomology and base change for proper flat finite presentation morphisms giving a perfect E = Rf\_*O\_X whose formation commutes with base change
\item the jump-locus lemma: where H\textasciicircum{}0 of a perfect complex is locally free of rank 1 the locus is universal closed and the map factors
\item the diagonal section argument splitting a direct image with a section
\item the cut-complex lemma splitting a perfect complex where H\textasciicircum{}0 fibre-surjects
\item the ring fact that an algebra invertible as a module over the base is the base
\end{itemize}\medskip
\ifdefined\hypertarget\hypertarget{sourcepointer-1}{}\fi
\noindent\textit{Source reference.} Original reference 0B91: sources/stacks-spaces-perfect/pinned-original-source.tex, source lines 6442--6442. Prerequisite pointer invoked at proof line(s) 6852; it names a standard fact of the same chapter that the counted argument uses, and every step of the counted proof built on it is retained verbatim in the delivered body.\par\medskip
\ifdefined\hypertarget\hypertarget{sourcepointer-2}{}\fi
\noindent\textit{Source reference.} Original reference 0BDI: sources/stacks-spaces-perfect/pinned-original-source.tex, source lines 6667--6667. Prerequisite pointer invoked at proof line(s) 6859; it names a standard fact of the same chapter that the counted argument uses, and every step of the counted proof built on it is retained verbatim in the delivered body.\par\medskip
\ifdefined\hypertarget\hypertarget{sourcepointer-3}{}\fi
\noindent\textit{Source reference.} Original reference 0BDL: sources/stacks-spaces-perfect/pinned-original-source.tex, source lines 6788--6788. Prerequisite pointer invoked at proof line(s) 6860; it names a standard fact of the same chapter that the counted argument uses, and every step of the counted proof built on it is retained verbatim in the delivered body.\par\medskip
\ifdefined\hypertarget\hypertarget{sourcepointer-4}{}\fi
\noindent\textit{Source reference.} Original lemma 0A1U: sources/stacks-chow/tags.txt, source lines 12682--12682. Prerequisite pointer invoked at proof line(s) 6891; it names a standard fact of a sibling chapter that the counted argument uses, and every step of the counted proof built on it is retained verbatim in the delivered body.\par\medskip
\begin{lemma}
\label{lemma-proper-flat-h0}
Let $S$ be a scheme.
Let $f : X \to Y$ be a morphism of algebraic spaces over $S$. Assume
\begin{enumerate}
\item $f$ is proper, flat, and of finite presentation, and
\item for a morphism $\Spec(k) \to Y$ where $k$ is a field, we have
$k = H^0(X_k, \mathcal{O}_{X_k})$.
\end{enumerate}
Then we have
\begin{enumerate}
\item[(a)] $f_*\mathcal{O}_X = \mathcal{O}_Y$ and
this holds after any base change,
\item[(b)] \'etale locally on $Y$ we have
$$
Rf_*\mathcal{O}_X = \mathcal{O}_Y \oplus P
$$
in $D(\mathcal{O}_Y)$
where $P$ is perfect of tor amplitude in $[1, \infty)$.
\end{enumerate}
\end{lemma}

\begin{proof}
It suffices to prove (a) and (b) \'etale locally on $Y$, thus we may and do
assume $Y$ is an affine scheme.
By cohomology and base change
(Lemma \ref{lemma-flat-proper-perfect-direct-image-general})
the complex $E = Rf_*\mathcal{O}_X$
is perfect and its formation commutes with arbitrary base change.
In particular, for $y \in Y$ we see that
$H^0(E \otimes^\mathbf{L} \kappa(y)) =
H^0(X_y, \mathcal{O}_{X_y}) = \kappa(y)$.
Thus $\beta_0(y) \leq 1$ for all $y \in Y$ with notation as in
Lemma \ref{lemma-jump-loci}. Apply
Lemma \ref{lemma-locally-closed-where-H0-locally-free}
with $a = 0$ and $r = 1$. We obtain a universal closed subscheme
$j : Z \to Y$ with $H^0(Lj^*E)$ invertible characterized
by the equivalence of (4)(a), (b), and (c) of the lemma.
Since formation of $E$ commutes with base change, we have
$$
Lf^*E = R\text{pr}_{1, *}\mathcal{O}_{X \times_Y X}
$$
The morphism $\text{pr}_1 : X \times_Y X$ has a section
namely the diagonal morphism $\Delta$ for $X$ over $Y$.
We obtain maps
$$
\mathcal{O}_X \longrightarrow R\text{pr}_{1, *}\mathcal{O}_{X \times_Y X}
\longrightarrow \mathcal{O}_X
$$
in $D(\mathcal{O}_X)$ whose composition is the identity. Thus
$R\text{pr}_{1, *}\mathcal{O}_{X \times_Y X} = \mathcal{O}_X \oplus E'$
in $D(\mathcal{O}_X)$. Thus $\mathcal{O}_X$ is a direct summand of
$H^0(Lf^*E)$ and we conclude that $X \to Y$ factors through $Z$
by the equivalence of (4)(c) and (4)(a) of the lemma cited above.
Since $\{X \to Y\}$ is an fppf covering, we have $Z = Y$.
Thus $f_*\mathcal{O}_X$ is an invertible $\mathcal{O}_Y$-module.
We conclude $\mathcal{O}_Y \to f_*\mathcal{O}_X$ is an isomorphism
because a ring map $A \to B$ such that $B$ is invertible as an $A$-module
is an isomorphism. Since the assumptions are preserved under base
change, we see that (a) is true.

\medskip\noindent
Proof of (b). Above we have seen that for every $y \in Y$ the map
$\mathcal{O}_Y \to H^0(E \otimes^\mathbf{L} \kappa(y))$ is surjective.
Thus we may apply
More on Algebra, Lemma \ref{more-algebra-lemma-better-cut-complex-in-two}
to see that in an open neighbourhood of $y$ we have
a decomposition $Rf_*\mathcal{O}_X = \mathcal{O}_Y \oplus P$
\end{proof}

\end{document}
